\section{Common tools for the later moment bounds}
\label{s:full-isotypic-section}
The every-occurrence theorem has supplied the full-deck mixing bound.
We now vary the information retained in a moment estimate. Some later
bounds optimize over signed factors and marks; others keep prescribed
hooks, minimize a deletion cost, or pair two diagram budgets. Their
optimization domains and moment exponents differ, even when they give
the same mixing order.

The shared geometric input is a row--column overlap at a fixed
placement of distinct marks. We first convert the compact-type density
comparison to full symmetric-group isotypes. Positive-mixture and
coefficient-integral formulas support the alternate density arguments;
the section ends with a marked recurrence at any hook parameter and
real Schatten order at least two. Each later section checks the local
hypotheses and budget inequalities needed to close its own induction.

\subsection{Guide to the later estimates}\label{s:later-routes}
Write \(Q_n=B_n^*B_n\) and \(\lvert B_n\rvert=Q_n^{1/2}\). All traces
in this guide are unnormalized; in particular,
\(\operatorname{Tr}Q_n^r=\operatorname{Tr}\lvert B_n\rvert^{2r}\).
The linked statements give the exact domains and constants.
Theorem~\ref{s:spin-moment} bounds
\(D_\lambda\operatorname{Tr}_{V_\lambda}Q_n^r\) for every prescribed
signed occurrence, retaining its concatenated entropy and clipped
remainder cost. The table compares the information retained by later
estimates and the arguments that prove them.

\begingroup
\renewcommand{\arraystretch}{1.05}
\begin{longtable}{@{}>{\raggedright\arraybackslash}p{0.22\linewidth}>{\raggedright\arraybackslash}p{0.40\linewidth}>{\raggedright\arraybackslash}p{0.32\linewidth}@{}}
\caption{Later estimates and their proof mechanisms.}\label{tab:s-later-routes}\\
\toprule
Route & Retained data and estimate & Proof mechanism\\
\midrule
\endfirsthead
\toprule
Route & Retained data and estimate & Proof mechanism\\
\midrule
\endhead
\bottomrule
\endfoot
Arbitrary-hook recurrence\newline Proposition~\ref{s:general-marked-recurrence} &
Every prescribed hook at any integer \(q\ge1\) is allowed. The recurrence retains real Schatten order \(p\ge2\), the full error, and a penalty for large child dimensions. &
Positive restriction to pointwise stabilizers combines with full symmetric-group isotypic overlap. No relation between the hook parameter and the child size is required.\\
\addlinespace[3pt]
Named objects\newline Proposition~\ref{ten:named:main} &
Every prescribed signed occurrence, with fixed numerical entropy and clipped remainder coefficients, is bounded at one fixed square moment. &
A direct specialization is followed by an independent crowded-cell and coefficient-integral proof.\\
\addlinespace[3pt]
Lowering and hook entropy\newline \eqref{ten:lowering:eq:1} and Proposition~\ref{ten:lowering:main} &
The squared operator norm is bounded for every \(2\le k\le n/2\). The trace bound minimizes Frobenius entropy plus a clipped remainder cost over permitted hook subdiagrams, with bounded real exponents of \(\lvert B_n\rvert\). &
Ordinary lowering and a pair count give the norm estimate; a distinct signed tensor interpolation gives the hook budget.\\
\addlinespace[3pt]
Minimized tags\newline Proposition~\ref{ten:compressed:main} &
Signed factors and a regular tag block are optimized. The budget can be negative and has an allowance; the overlap retains ranks inside local Specht modules. &
Sparse path encounters give the small-defect norm bound. Interpolating the covariance and multiplicity bounds retains the local ranks.\\
\addlinespace[3pt]
Simultaneous bounds\newline Proposition~\ref{ten:simultaneous:main} &
One fixed integer square moment has both a dimension bound and a marked bound for every permitted hook whose marked complement meets the stated density threshold. &
A specialized recurrence also has an independent proof through powers-of-two trace inequalities and weight-space overlap. Harmonic-tuple coupling and a dimension gain for shrinking hooks close the two assertions together.\\
\addlinespace[3pt]
Inverse-density budget\newline Proposition~\ref{s:inverse-bounded-exponent} &
A minimum over permitted hooks retains the unclipped mark-entropy cost at large levels, with bounded real Schatten exponents. Sparse levels have a separate bound. &
The marked-grid estimate (Proposition~\ref{ten:inverse:main}) has an ordered-density proof. Allowances extend each child estimate to every inherited parent hook before the exponent recursion closes.\\
\addlinespace[3pt]
Deletion budget\newline Theorem~\ref{rec:deletion-moment} &
For every \(n=2^d\) with \(d\ge1\) and every partition, a clipped minimum over all subdiagrams bounds the dimension-weighted trace, with real square-moment orders between one and a fixed finite bound. &
A near-minimizer in a thin hook and signed overlap handle large sizes; norm gaps in finitely many blocks supply one base before the exact recursion.\\
\addlinespace[3pt]
Paired budgets\newline Theorem~\ref{rec:hybrid-main} &
For nontrivial diagrams of height at most \(n/2\), two bounds for the same fixed integer square moment pay, respectively, tuple-coordinate losses and long row and column chains with distinct markers. &
Separate tuple and marked tensor estimates are joined by a diagram comparison, with an explicit finite base.\\
\end{longtable}
\endgroup

\subsection{Full symmetric-group isotypes}
The conversion from compact types to full symmetric-group isotypes is
static and uses no moment estimate. Here
$V=\mathbb C^q\oplus\mathbb C^q$ again has one even and one odd
summand, for any integer $q\ge1$. A full symmetric-group isotype
collects the copies of its Specht module from every compatible compact
color type. It includes all those multiplicity copies, whereas
$Q_{a,b}$ fixes one compact type.

\begin{theorem}[Signed density domination for a full isotypic projection]
\label{thm:signed-isotypic-density}
Let $P_\rho$ be the symmetric-group isotypic projection of shape
$\rho\vdash p$ on the signed tensor power $V^{\otimes p}$.
There is a probability distribution on even densities such that
\begin{equation}
 P_\rho\preceq D_\rho B^{10q^2}\mathbb E R^{\otimes p},
 \qquad B=N+q+1,\quad N\ge p.
 \label{eq:signed-isotypic-density}
\end{equation}
If $P_\rho=0$, the assertion is read with any density distribution.
\end{theorem}
\begin{proof}
Let $\mathcal A_\rho$ consist of the pairs $(a,b)$ for which
$V_\rho$ occurs in $\mathcal M_{a,b}$. Orthogonal sector
decomposition gives $P_\rho\preceq\sum_{\mathcal A_\rho}Q_{a,b}$.
Apply Lemma~\ref{lem:signed-type-density} and combine its probability
mixtures, with weights proportional to $e^{F(a,b)}u(a,b)$.
Their total coefficient is at most
\[
 \sum_{\mathcal A_\rho}e^{F(a,b)}u(a,b)
 \le D_\rho B^{7q^2}B^{2q^2}B^{q^2}
 =D_\rho B^{10q^2},
\]
by Lemma~\ref{lem:signed-hook-carriers}.
\end{proof}

\begin{corollary}[Full-isotypic angle on a punctured rectangle]
\label{cor:signed-isotypic-angle}
On the board $W$ of Theorem~\ref{thm:one-sided-type-angle}, let
$P_\gamma$ be a full $S_p$-isotypic projection, and let $P_R,P_C$
be products of full symmetric-group isotypic projections on the
occupied row and column sites. Write their diagrams as $\rho_i$ and
$\sigma_j$, and set $D_R=\prod_iD_{\rho_i}$,
$D_C=\prod_jD_{\sigma_j}$. With $N=sm$ and $B=N+q+1$,
\begin{equation}
 \|P_C P_\gamma P_R\|_{\rm op}^2
 \le\min\left\{1,
  \frac{D_RD_C}{D_\gamma}B^{10(s+m)q^2}e^l\right\}.
 \label{eq:signed-isotypic-angle}
\end{equation}
Empty blocks use the trivial partition and dimension one. The
conclusion is uniform over the location of the $l$ missing cells.
\end{corollary}
\begin{proof}
Repeat the proof of Theorem~\ref{thm:one-sided-type-angle}, now using
Theorem~\ref{thm:signed-isotypic-density} for each local projector.
The only changed global estimate is
\[
 \|T P_\gamma\|_{\rm op}\le D_\gamma^{-1}
\]
for the identical-factor density $T=\bar R^{\otimes W}$.
Indeed $T$ has trace one and commutes with the signed permutation
action. On the $\gamma$ isotypic space it has the form
$I_{V_\gamma}\otimes T_\gamma$, with $T_\gamma\succeq0$, so
$D_\gamma\|T_\gamma\|\le
D_\gamma\operatorname{Tr}T_\gamma\le1$.
The full-group projection commutes with both subgroup projections,
and $T$ commutes with all of them. The support and missing-cell
arguments are otherwise identical. Taking adjoints identifies the
two displayed orders of projections.
\end{proof}


\begin{proposition}[One-sided overlap for positive marked factors]
\label{s:one-sided-marked-overlap}
Let $W$ contain $n-t$ occupied cells in an $a$ by $b$ board, $n=ab$.
Assume $\omega\vdash n-t$ and every prescribed row and column type
$\eta_g$ are $(q,q)$-hook shapes: their $(q+1)$st part is at most $q$.
Here $1\le q\le n$.
Suppose $X=\prod_{g\text{ row}}X_g$ and
$Y=\prod_{g\text{ column}}Y_g$ are positive group-algebra elements,
with each $X_g,Y_g$ supported on the indicated local type. Write
$c_g$ for its unnormalized trace on one copy of $V_{\eta_g}$ and
$S=\sum_g\log D_{\eta_g}$. Then
\[
 \operatorname{Tr}_{\omega}(XY)
 \le\Bigl(\prod_g c_g\Bigr)
 \exp\{C(a+b+1)q^2\log(n+1)+t
                  +\min(0,S-\log D_\omega)\}.
\]
The constant is absolute. On the $\omega$-isotypic part of the signed
tensor power, the same bound is multiplied by its multiplicity
$w_\omega$.
\end{proposition}
\begin{proof}
In the signed tensor space the global symmetric-group projector
$P_\omega$ commutes with $X,Y$. Their local supporting projections
are $P_R,P_C$. Spectral decomposition of their positive compressions gives
\[
 \operatorname{Tr}(P_\omega XY)
 \le \|P_CP_\omega P_R\|^2\operatorname{Tr}X\operatorname{Tr}Y.
\]
Indeed each squared inner product of a row eigenvector and a column
eigenvector is bounded by the squared overlap, and the sum of their
eigenvalue products is the product of the two traces. Each local
trace is $w_{\eta_g}c_g$, with
$w_{\eta_g}\le(n+q+1)^{10q^2}$ by
Lemma~\ref{lem:signed-hook-carriers}. Apply
Corollary~\ref{cor:signed-isotypic-angle}; it bounds the squared overlap
by both $1$ and
$\exp\{S-\log D_\omega+t+C(a+b)q^2\log(n+1)\}$.
Finally $\operatorname{Tr}(P_\omega XY)=w_\omega
\operatorname{Tr}_\omega(XY)$ and $w_\omega\ge1$.
Taking the smaller bound and absorbing all carrier powers proves the
claim, including empty lines and the all-hole board.
\end{proof}

\subsection{Positive factorizations and missing cells}\label{ten:factorizations}
\begin{lemma}[Positive domination by mixtures]\label{ten:positive-mixtures}
Let $A,B\succeq0$ and suppose $A\preceq c_A\int A_x\,d\mu(x)$ and $B\preceq c_B\int B_y\,d\nu(y)$, where $c_A,c_B\ge0$, all integrands are positive, and $\mu,\nu$ are probability measures. Then, for every $P$,
\[
 \|A^{1/2}PB^{1/2}\|\le\sqrt{c_Ac_B}\sup_{x,y}\|A_x^{1/2}PB_y^{1/2}\|.
\]
\end{lemma}
\begin{proof}
Define $Fv=(A_x^{1/2}v)_x$ and $Gv=(B_y^{1/2}v)_y$. Positive domination gives contraction factorizations $A^{1/2}=\sqrt{c_A}UF$ and $B^{1/2}=\sqrt{c_B}G^*V^*$: for the first define $U(Fv)=A^{1/2}v/\sqrt{c_A}$ on the range, and use its adjoint for the second. Zero constants give the assertion directly. The integral operator $FPG^*$ has the displayed one-pair operators as its kernel. Its norm is bounded by their supremum by Cauchy--Schwarz and the probability normalizations. Multiplication by the two contractions proves the claim.
\end{proof}

\begin{lemma}[Scalar missing-cell estimate]\label{ten:missing-cells}
Let $\rho_i$ and $\sigma_j$ be trace-one positive matrices indexed by $a$ rows and $b$ columns. Put $\bar\rho=a^{-1}\sum_i\rho_i$ and use its inverse on its support. For a unitary $U$ set
\[
 z_{ij}=\operatorname{Tr}(\bar\rho^{-1/2}\rho_i\bar\rho^{-1/2}U\sigma_jU^*)\ge0.
\]
Then $\sum_{i,j}z_{ij}\le ab$. On any $ab-t$ occupied cells their product is at most $e^t$.
\end{lemma}
\begin{proof}
Summing the normalized row matrices gives $a$ times the support projection of $\bar\rho$. Since every $\sigma_j$ has trace one, the full sum is at most $ab$. Arithmetic--geometric means give $(ab/(ab-t))^{ab-t}\le e^t$; the empty product is one.
\end{proof}

\subsection{Extracting one compact type by matrix coefficients}

The alternate density arguments below use a global compact-type
projection between two products of local densities. The following
matrix-coefficient formula turns that middle operator into an integral
of tensor actions, with its carrier dimension explicit.

\begin{lemma}[A coefficient integral for one compact type]
\label{s:compact-coefficient-extraction}
Let $K$ be a compact group, let $\rho$ be a finite-dimensional unitary
representation on $\mathcal H$, and let $\pi$ be an irreducible
unitary representation on $\mathcal U$ of dimension $M$.
Identify the $\pi$-isotypic subspace with
$\mathcal U\otimes\mathcal M$. For $A\in\operatorname{End}(\mathcal U)$,
let $\widetilde A$ act there as $A\otimes I_{\mathcal M}$ and be
zero on all other isotypes. With normalized Haar measure,
\[
 \widetilde A=\int_K f_A(g)\rho(g)\,dg,
 \qquad f_A(g)=M\operatorname{Tr}(A\pi(g)^*),
\]
and
\[
 \int_K|f_A(g)|\,dg
 \le\sqrt M\,\|A\|_{\mathrm{HS}}
 \le M\|A\|_{\mathrm{op}}.
\]
If the type does not occur, $\widetilde A=0$ and the same identity
holds.
\end{lemma}
\begin{proof}
Apply compact matrix-coefficient orthogonality
\cite[Theorem IV.3.8(i)--(ii)]{Morel2018} in an orthonormal basis of
$\mathcal U$. On each copy of $\pi$, the integral has matrix $A$;
on every inequivalent irreducible it is zero. The same orthogonality
gives
$\int_K|f_A|^2=M\operatorname{Tr}(A^*A)$.
Cauchy--Schwarz for the probability Haar measure and
$\|A\|_{\mathrm{HS}}\le\sqrt M\|A\|_{\mathrm{op}}$
prove the bounds.
\end{proof}

\subsection{A marked recurrence with the full error term}
\label{s:general-marked-recurrence-section}

The full-isotypic overlap now gives a recurrence for a prescribed
marked complement. The estimate retains the hook parameter and the
real Schatten exponent.
This is useful when a subdiagram permitted at a parent scale is larger than
those permitted in the child induction.  In particular, no relation between
the hook parameter and the child size is imposed.

Let $B_m$ be one coordinate sweep on $m$ positions, for $m$ a power of two,
and put
\[
 H_\xi=\log D_\xi,\qquad
 E_m(\xi;p)=\log\bigl(D_\xi\operatorname{Tr}_\xi|B_m|^p\bigr),
 \qquad j_m(l)=l\log(m/l).
\]
We set $j_m(0)=0$, $H_\varnothing=0$, and $E_m(\xi;p)=-\infty$ when
the trace vanishes.  All traces are unnormalized.  A partition $\eta$ lies
in the $q$-hook if $\eta_{q+1}\le q$.  Thus this definition includes the
empty partition.

\begin{proposition}[Marked recurrence at an arbitrary hook parameter]
\label{s:general-marked-recurrence}
Let $n=ab\ge4$, where $a,b\ge2$ are powers of two and
$|\log_2a-\log_2b|\le1$.  Arrange the positions as $a$ rows of length
$b$, so that $B_n=CR$, where $R$ consists of the $a$ independent row
sweeps $B_b$ and $C$ consists of the $b$ independent column sweeps $B_a$.
Let $p\ge2$ be real and let $q\ge1$ be any integer.  For every
$\lambda\vdash n$ with positive moment and every $q$-hook subpartition
$\omega\subseteq\lambda$, put $t=n-|\omega|$.  Then
\begin{align}
 E_n(\lambda;p)+j_n(t)
 &\le \max_{\mathcal A}\left\{
    \sum_g\bigl(E_{m_g}(\xi_g;p)+j_{m_g}(l_g)\bigr)
      -\left(\sum_g H_{\eta_g}-H_\omega\right)_+
                         \right\}\notag\\
 &\hspace{12mm}
    +C\left[t+(a+b)(q^2+n^{1/4})\log(n+1)\right].
 \label{s:eq:general-marked-recurrence}
\end{align}
Here $g$ ranges over all rows and columns, and $m_g=b$ on a row and
$m_g=a$ on a column.  The set $\mathcal A$ consists of choices satisfying
\[
 0\le l_g\le m_g,\qquad
 \sum_{g\text{ row}}l_g=\sum_{g\text{ column}}l_g=t,
\]
\[
 \xi_g\vdash m_g,\quad E_{m_g}(\xi_g;p)>-\infty,
 \qquad \eta_g\subseteq\xi_g,\quad |\eta_g|=m_g-l_g,
 \quad (\eta_g)_{q+1}\le q.
\]
The set $\mathcal A$ is nonempty when the parent moment is positive.
The constant $C$ is absolute, independent of $p,q,a,b,\lambda,\omega$.
\end{proposition}

\begin{proof}
We can replace $q$ by $Q=\min(q,n)$ in the proof: every partition of
size at most $n$ has the same $q$-hook and $Q$-hook status.  This permits
the use of the signed tensor space with even and odd dimensions $Q$
without imposing an upper bound on the stated parameter $q$.

First compare the parent moment with positive row and column factors.
Set
\[
 X=(RR^*)^{p/2},\qquad Y=(C^*C)^{p/2}.
\]
Araki--Lieb--Thirring at exponent $p/2\ge1$, together with the polar
decomposition of $R$, gives
\[
 \operatorname{Tr}_\lambda|CR|^p
 \le \operatorname{Tr}_\lambda XY.
\]
This is the only step involving the value of $p$.  The two positive
operators factor over the rows and the columns respectively.  Expanding
each local factor by its central symmetric-group isotypic projections
gives positive factors supported on types $\xi_g$, with irreducible
trace $\operatorname{Tr}_{\xi_g}|B_{m_g}|^p$; replacing $B_{m_g}$ by its
adjoint does not change that trace.

Fix one such choice of the $\xi_g$.  Let
$V=\mathbb C^Q\oplus\mathbb C^Q$, with the first summand even and the
second odd, and let $\Pi_\omega$ be the $\omega$-isotypic projector on
the signed tensor power $V^{\otimes(n-t)}$.  Its multiplicity $w_\omega$
is positive.  Add $t$ distinct even marks, each used once.  The direct
sum over their ordered placements of the spaces
$\Pi_\omega V^{\otimes(n-t)}$ is an $S_n$-module $\mathcal K$.  More
precisely,
\[
 \mathcal K\cong
 \operatorname{Ind}_{S_{n-t}}^{S_n}
       \bigl(V_\omega\otimes\mathbb C^{w_\omega}\bigr).
\]
Branching shows that $\lambda$ occurs with multiplicity
$w_\omega f^{\lambda/\omega}$.  Each irreducible trace of $XY$ is
nonnegative because $X,Y$ are positive.  Consequently
\[
 \operatorname{Tr}_\lambda XY
 \le \frac{\operatorname{Tr}_{\mathcal K}XY}
              {w_\omega f^{\lambda/\omega}}.
 \tag{$\ast$}
\]

Insert the ordered-placement projections between $X$ and $Y$ in the
numerator.  Only a common diagonal placement contributes: a row
permutation preserves each mark's row and a column permutation preserves
its column, so an intermediate placement contributing to a return must
agree with the initial placement in both coordinates
(Figure~\ref{fig:marked-grid}).  Fix such a
placement $z$, and write $l_g$ for its mark counts.

On a block of size $m=m_g$, diagonal compression keeps precisely the
coefficients of the subgroup fixing the $l=l_g$ marked sites pointwise.
Coefficient restriction to $S_{m-l}$ preserves positivity, by compression
of the regular representation to that subgroup.  Its regular trace is
\[
 \frac{D_{\xi_g}\operatorname{Tr}_{\xi_g}|B_m|^p}{(m)_l},
 \qquad (m)_l=m(m-1)\cdots(m-l+1).
\]
Indeed the identity coefficient is unchanged, and the two regular
traces multiply that coefficient by $m!$ and $(m-l)!$ respectively.
The restricted element is supported only on $\eta_g\subseteq\xi_g$:
if a subgroup type is absent from the restriction of $V_{\xi_g}$, its
central projector annihilates the original element, and coefficient
restriction commutes with multiplication by subgroup elements.  After
splitting by these subgroup types, the trace of a local positive factor
on $V_{\eta_g}$ is therefore at most
\[
 c_g=\frac{\exp E_{m_g}(\xi_g;p)}{(m_g)_{l_g}D_{\eta_g}}.
 \tag{$\ast\ast$}
\]
Only $Q$-hook types occur in the remaining signed tensor factors.

The one-sided signed-tensor overlap estimate
(Proposition~\ref{s:one-sided-marked-overlap}) applies to these positive
local factors.  With $S=\sum_g H_{\eta_g}$, it bounds the trace for this
placement and these types by
\[
 w_\omega\Bigl(\prod_g c_g\Bigr)
 \exp\!\left(C\bigl((a+b+1)Q^2\log(n+1)+t\bigr)
                +\min(0,S-H_\omega)\right).
\]
The placement compression acts in the core permutation algebra, so
$\Pi_\omega$ commutes with each compressed factor; this is precisely
the projection placement required by that overlap estimate.  Substituting
$(\ast\ast)$, the logarithm of the last display is at most
\[
 \log w_\omega+\sum_g E_{m_g}(\xi_g;p)
       -\sum_g\log(m_g)_{l_g}-\max(S,H_\omega)
       +C\bigl((a+b+1)Q^2\log(n+1)+t\bigr).
 \tag{$\ast\ast\ast$}
\]

It remains to count the ordered marks.  For fixed row and column count
vectors, an ordered placement determines an assignment of each mark to
a row and to a column.  Hence its number is at most
\[
 \frac{t!}{\prod_{g\text{ row}}l_g!}
 \frac{t!}{\prod_{g\text{ column}}l_g!}.
\]
There is no additional choice of mark-label allocations.  The entropy
bound for each multinomial and $(m)_l\ge(m/e)^l$ show that the logarithm
of this number, minus $\sum_g\log(m_g)_{l_g}$, is at most
\begin{align*}
 2t\log t-\sum_g l_g\log l_g-t\log n+2t
   &=\sum_g j_{m_g}(l_g)-2j_n(t)+2t.
\end{align*}
All expressions are interpreted continuously at zero.  Finally, the
number of count and type choices has logarithm at most
\[
 C\bigl(a\sqrt b+b\sqrt a+a+b\bigr)\log(n+1)
 \le C(a+b)n^{1/4}\log(n+1).
\]
For example, a partition of size at most $m$ is specified by its Durfee
square and at most $2\sqrt m$ row and column lengths; the hole count
adds one integer.  This count also covers the preceding total-type
expansion.

By assigning entries of a standard tableau separately to $\omega$ and
$\lambda/\omega$ and discarding the boundary inequalities,
\[
 D_\lambda\le\binom ntD_\omega f^{\lambda/\omega},
 \qquad \log\binom nt\le j_n(t)+t.
\]
Use these inequalities in $(\ast)$, sum the bounds $(\ast\ast\ast)$,
and add $j_n(t)$ to the parent logarithm.  The multiplicity
$w_\omega f^{\lambda/\omega}$ cancels, the mark-entropy terms cancel,
and the remaining dimension term is
\[
 H_\omega-\max(S,H_\omega)=-(S-H_\omega)_+.
\]
This proves the recurrence.  If every admissible term vanished, the
positive parent moment would also vanish by these inequalities; hence
the stated maximum is over a nonempty set.
\end{proof}

\begin{corollary}[Exact specialization to the inverse-density budget]
\label{s:inverse-budget-specialization}
Under the hypotheses of Proposition~\ref{s:general-marked-recurrence},
let $0\le\theta\le1$, let $c$ be real, and let $\mathcal B_a,\mathcal B_b\ge0$.
Suppose that, for $m\in\{a,b\}$, every $\xi\vdash m$ and every
$q$-hook $\eta\subseteq\xi$, with $l=m-|\eta|$, satisfy
\begin{equation}
 E_m(\xi;p)\le \mathcal B_m+\theta H_\eta+c\,l\log m-j_m(l).
 \label{s:eq:inverse-child-budget}
\end{equation}
Then every prescribed $q$-hook $\omega\subseteq\lambda$, with
$t=n-|\omega|$, satisfies
\begin{align}
 E_n(\lambda;p)
 &\le\theta H_\omega+c\,t\log n-j_n(t)+a\mathcal B_b+b\mathcal B_a\notag\\
 &\hspace{9mm}+C\left[t+(a+b)(q^2+n^{1/4})\log(n+1)\right].
 \label{s:eq:inverse-budget-specialization}
\end{align}
The constant $C$ does not depend on $p,q,\theta,c,\mathcal B_a,\mathcal B_b$.
\end{corollary}
\begin{proof}
For an admissible choice in the recurrence put $S=\sum_g H_{\eta_g}$.
The sum of the child allowances is $a\mathcal B_b+b\mathcal B_a$, and
\[
 \sum_g l_g\log m_g=t\log b+t\log a=t\log n.
\]
Thus the expression being maximized is at most
\[
 a\mathcal B_b+b\mathcal B_a+c\,t\log n+\theta S-(S-H_\omega)_+.
\]
For $S,H_\omega\ge0$ and $0\le\theta\le1$,
\[
 \theta S-(S-H_\omega)_+\le\theta H_\omega:
\]
if $S\le H_\omega$ this is monotonicity; otherwise its left side equals
$\theta H_\omega-(1-\theta)(S-H_\omega)$.  Subtract $j_n(t)$.
\end{proof}

In particular, taking $\theta=1/2$, $c=1/100$ and
$q\le n^{1/16}$ gives the inverse-grid estimate with its full stated
error and its additive child allowances. Section~\ref{ten:inverse}
also derives this specialization through its ordered density product.

The common recurrence exposes the child-dimension penalty and mark
entropy before a budget is chosen. It does not by itself supply the
sparse estimate or the scale comparison needed to close an induction.
We begin the later bounds with named objects, whose budget still
prescribes a signed occurrence. Its direct deduction from the first
theorem fixes the scope; the rest of that section develops a second
proof through crowded cells and coefficient integrals.
